\documentclass[a4,12pt]{article}
\usepackage[utf8]{inputenc}
\usepackage[T2A]{fontenc}
\usepackage[english, russian]{babel}


\usepackage{csquotes}
\usepackage{epigraph}
\usepackage{tcolorbox}
% Определяем окружение для Формата ввода
\newtcolorbox{inputformat}[1][]{colback=blue!5!white, colframe=blue!65!black,
    fonttitle=\bfseries, title=Формат ввода, #1}

% Определяем окружение для Формата вывода
\newtcolorbox{outputformat}[1][]{colback=green!5!white, colframe=green!65!black,
    fonttitle=\bfseries, title=Формат вывода, #1}

% Определяем окружение для комментария
\newtcolorbox{exercisecomment}[1][]{colback=yellow!5!white, colframe=yellow!80!black, fonttitle=\bfseries, title=Комментарий, #1}

% Определяем окружение для замечания
\newtcolorbox{exercisenote}[1][]{colback=red!5!white, colframe=red!80!black, fonttitle=\bfseries, title=Замечание, #1}


\usepackage{geometry}
%\geometry{papersize={20.3 cm,25.4 cm}}
\geometry{left=2cm}
\geometry{right=2cm}
\geometry{top=2cm}
\geometry{bottom=2cm}


\usepackage{amsmath, amsfonts, amsthm, amssymb, amsopn, amscd}
\usepackage{enumerate}
\usepackage{enumitem}
\usepackage[mathscr]{eucal}

\usepackage{hyperref}
\hypersetup{unicode=true,final=true,colorlinks=true}

\theoremstyle{remark}
%\newtheorem{exercises}{Упражнения}
\newtheorem{exercise}{\textbf{Упражнение}}[section]
\renewcommand{\theexercise}{\textbf{\arabic{exercise}}}
%\renewcommand{\theexercise}{\textbf{\#\arabic{exercise}}}


\title{Листок 06. Методы работы со словарями в Python}

\author{А.Н. Лата}

\begin{document} 

%\maketitle

\section*{\centering Листок 06. Методы работы со словарями в Python}

\begin{exercisenote}[title=Замечания]
\begin{itemize}
    \item для решение приведенных ниже упражнений не требуется создавать (определять) пользовательские функции, циклы и условные операторы.
    \item при решении упражнений полученные результаты выведите на экран с помощью функции {\color{blue}{print()}}.
    \item позаботесь о том, чтобы выводимый на экран результат был снабжен информацией о нем, там где это необходимо.
    \item не забывайте писать комментарии к вашему коду
\end{itemize}
\end{exercisenote}

\textbf{Упражнения}

\bigskip
\noindent\textbf{Часть~I. Создание словарей}

\begin{enumerate}

    \item Создайте пустой словарь и выведите его. Выведите также тип созданного объекта с помощью функции \textbf{type()}.

    \item Создайте словарь \texttt{товар}, содержащий три пары: название товара, его цена и количество на складе. Выведите словарь.

    \item Создайте словарь \texttt{сотрудник} с помощью конструктора \textbf{dict()}, указав ключи: имя, должность, оклад. Выведите словарь.

    \item Даны два списка: \texttt{['выручка', 'затраты', 'прибыль']} и \texttt{[1\,500\,000, 900\,000, 600\,000]}. Создайте словарь из этих списков с помощью \textbf{dict(zip(...))}. Выведите результат.

    \item Создайте словарь из списка кортежей \texttt{[('USD', 89.5), ('EUR', 97.3), ('CNY', 12.4)]} с помощью функции \textbf{dict()}. Выведите словарь.

    \item Создайте словарь, в котором ключами являются элементы списка \texttt{['январь', 'февраль', 'март']}, а все значения равны нулю. Используйте метод \textbf{dict.fromkeys()}. Выведите результат.

    \item Создайте словарь с данными о компании: название, год основания, количество сотрудников, выручка. Выведите количество пар ключ--значение с помощью функции \textbf{len()}.

    \item Даны кортеж городов \texttt{('Москва', 'СПб', 'Казань')} и список объёмов продаж \texttt{[850\,000, 420\,000, 310\,000]}. Создайте словарь с помощью \textbf{dict(zip(...))}. Выведите словарь.

\end{enumerate}

\bigskip
\noindent\textbf{Часть~II. Доступ к значениям и проверка ключей}

\begin{enumerate}[resume]

    \item Создайте словарь с курсами валют: USD, EUR, GBP. Обратитесь к курсу доллара и евро по ключу и вычислите их среднее значение. Выведите результат с двумя знаками после запятой.

    \item Создайте словарь с зарплатами трёх сотрудников. Обратитесь к зарплатам двух конкретных сотрудников по ключу и выведите их разницу.

    \item Создайте словарь остатков товаров на складе. Используйте метод \textbf{.get()} для обращения к товару, который \textbf{есть} в словаре. Выведите полученное значение.

    \item Создайте словарь остатков товаров на складе. Используйте метод \textbf{.get()} для обращения к товару, которого \textbf{нет} в словаре. В качестве значения по умолчанию передайте строку \texttt{'Нет в наличии'}. Выведите результат.

    \item Создайте словарь налоговых ставок с ключами \texttt{'НДС'}, \texttt{'НДФЛ'}, \texttt{'налог\_на\_прибыль'}. Проверьте, есть ли в словаре ключ \texttt{'НДС'}, с помощью оператора \textbf{in}. Выведите результат проверки.

    \item Используя словарь из упражнения~13, проверьте оператором \textbf{not~in}, отсутствует ли в словаре ключ \texttt{'акциз'}. Выведите результат проверки.

    \item Создайте словарь прайс-листа из пяти товаров. Выведите одной строкой название и цену второго товара, используя обращение по ключу и f-строку.

\end{enumerate}

\bigskip
\noindent\textbf{Часть~III. Изменение словаря}

\begin{enumerate}[resume]

    \item Создайте словарь \texttt{бюджет} с тремя статьями расходов. Добавьте в него новую статью \texttt{'реклама'} со значением~50\,000. Выведите обновлённый словарь.

    \item Создайте словарь с окладом сотрудника. Увеличьте оклад на $20\%$, перезаписав значение по ключу. Выведите словарь до и после изменения.

    \item Создайте словарь с ценами пяти товаров. Удалите одну запись с помощью оператора \textbf{del}. Выведите словарь до и после удаления.

    \item Создайте словарь с данными клиента: имя, телефон, баланс счёта. Удалите ключ \texttt{'телефон'} методом \textbf{.pop()} и сохраните удалённое значение в переменную. Выведите удалённое значение и обновлённый словарь.

    \item Создайте два словаря: \texttt{план} и \texttt{факт} с объёмами продаж по трём регионам. Обновите словарь \texttt{план} данными из словаря \texttt{факт} с помощью метода \textbf{.update()}. Выведите результат. \textbf{Обратите внимание}: что происходит со значениями при совпадении ключей?

    \item Создайте словарь ежемесячных продаж за три месяца. Используйте метод \textbf{.setdefault()}, чтобы добавить четвёртый месяц со значением по умолчанию~0. Вызовите \textbf{.setdefault()} повторно для \textbf{уже существующего} ключа и убедитесь, что значение не изменилось. Выведите словарь.

\end{enumerate}

\bigskip
\noindent\textbf{Часть~IV. Методы просмотра: .keys(), .values(), .items()}

\begin{enumerate}[resume]

    \item Создайте словарь с показателями финансовой отчётности компании: выручка, себестоимость, валовая прибыль, EBITDA. Выведите все ключи словаря, преобразовав их в список с помощью \textbf{list()}.

    \item Используя словарь из упражнения~22, выведите все значения словаря, преобразовав их в список с помощью \textbf{list()}.

    \item Используя словарь из упражнения~22, выведите все пары ключ--значение, преобразовав результат метода \textbf{.items()} в список. Обратите внимание на тип каждого элемента полученного списка.

    \item Создайте словарь с объёмами продаж пяти менеджеров. Преобразуйте \textbf{.items()} в список кортежей и выведите первый и последний элемент этого списка по индексу.

    \item Создайте словарь с ценами шести товаров. С помощью \textbf{list()} и \textbf{.keys()} получите список названий товаров. Выведите срез, содержащий первые три названия.

\end{enumerate}

\bigskip
\noindent\textbf{Часть~V. Агрегирующие функции}

\begin{enumerate}[resume]

    \item Создайте словарь с квартальной выручкой компании: Q1, Q2, Q3, Q4. Найдите годовую выручку с помощью \textbf{sum()}, применив её к значениям словаря. Выведите результат.

    \item Создайте словарь с объёмами продаж по регионам. С помощью \textbf{max()} найдите наибольший объём продаж. С помощью \textbf{min()} найдите наименьший. Выведите оба значения и их разницу.

    \item Создайте словарь с зарплатами пяти сотрудников. Вычислите среднюю зарплату, используя \textbf{sum()} и \textbf{len()}. Выведите результат.

    \item Создайте словарь с ценами товаров. С помощью \textbf{sorted()} и \textbf{.values()} получите список цен, отсортированный по возрастанию. Выведите результат. Повторите с параметром \textbf{reverse=True}.

    \item Создайте словарь с объёмами продаж менеджеров. С помощью \textbf{sorted()} и \textbf{.items()} получите список пар, отсортированных по объёму продаж (по убыванию). Используйте параметр \textbf{key=lambda x: x[1]}. Выведите первые три элемента списка (срез \texttt{[:3]}).

    \item Создайте словарь с курсами трёх валют. С помощью \textbf{sorted()} и \textbf{.keys()} выведите список валют, отсортированных по алфавиту.

\end{enumerate}

\end{document}